\section{A bounded restart race between certified interval sweeps}
\label{sec:orbit-race}

The width selector of Section~\ref{sec:orbit-direction} chooses an order
without observing the ensuing search. Its counterexample shows that
fewer cycles need not mean less work. The new optional
\code{sweep\_direction='race'} instead interrupts attempts at a measured
checkpoint allowance, tries the other direction and increases the allowance
geometrically. This gives a precise bound for the scheduling policy.
The unit of this bound is a cooperative checkpoint, not an elementary
bit operation or elapsed time. That distinction is essential because some
uninterrupted blocks sort or scan a whole collection.

The race is available in both \code{count\_orbits} and
\code{normal\_surface\_topology}. Forward remains the default; the
previous forward, reverse and width-selected paths retain their complete
records exactly. The race is a new search policy using the existing proof
formats and unchanged independent checkers.

\subsection{The restart policy and its exact accounting}

Retain a validated list of the $k$ input pairings on a universe of size
$M$, and let $b=\operatorname{bit\_length}(M)$. Set the initial allowance to
\[
 L=8(k+1)(b+1).
\]
The earlier width selector chooses which direction is tried first. At
round $r=0,1,\ldots$, both directions, in that fixed order, receive allowance
$L2^r$. Each attempt starts again from the original immutable pairings.
On every checkpoint the caller's cancellation callback runs first; then
the scheduler increments the attempt and total work counters. The attempt
is interrupted if its counter exceeds the allowance. Consequently, an
exhausted attempt costs exactly one more checkpoint than its allowance.
A completed attempt immediately returns its answer and ordinary certificate.

The interruption uses a fresh local exception class, so unrelated caller
exceptions propagate. Failed attempt frames and their partial certificates
are discarded. No intermediate orbit count is promoted to a result. The
winning certificate is exactly the one the corresponding fixed direction
would have emitted in an unlimited run. Thus independent verification
does not need to trust the direction selector, restart decisions, abandoned
states or scheduling counters.

The existing \code{max\_cycles} allowance is shared across every attempt.
A private hook records each begun AHT cycle, including one interrupted
halfway through. Each new attempt receives only the remaining cycle
allowance. If it is spent, the race returns incomplete with neither a
count nor a certificate. In a normal-surface query, the same accounting
continues across all surface, double, boundary and optional cone queries.
An empty universe can still finish under a zero-cycle allowance.

Public \code{cycles} reports all begun cycles. Ordinary structural statistics
describe the winning attempt, or the last attempt on an incomplete return.
An interrupted helper can leave those partial structural counters behind
its actual work; they are not used for the scheduling bound.
Additional \code{race\_*} counters separately report attempts, restarts,
work exhaustions, initial and final allowances, directional work and cycles,
winning work and cycles, and startup work. A final exhaustion at the cycle
cap is a work exhaustion even though it causes no further restart. Proof
event allowances count only the returned certificate, which contains no
abandoned prefixes. These different counters are intentionally not interchangeable.

\subsection{A competitive checkpoint bound}

Fix the input, merger scheduler, merger threshold and certificate-recording
setting. Let $W_0,W_1$ be the complete checkpoint counts of the two unlimited
fixed-direction runs, excluding the initial width-selection scan, and put
$W=\min(W_0,W_1)$. Assume no external cancellation or cycle cap stops the
computation. Let $B=L2^r$ be the least allowance at least $W$.
No preceding round can finish either direction, and some direction must
finish at allowance $B$. The preceding rounds spend at most
\[
 2\sum_{j=0}^{r-1}(L2^j+1)=2(B-L)+2r
\]
checkpoints. The final round spends at most $2B+1$: one exhausted attempt
and one completed attempt, or a single completed attempt. Therefore the
total attempt work $A$ satisfies
\[
 A\leq4B-2L+2r+1\leq12\max(L,W).
\]
The last inequality is deliberately coarse. For $r=0$ the first bound is
$2L+1$; for $r>0$, use $B<2W$ and $2r+1\leq2B$.
The initial selection contributes exactly $k+1$ additional checkpoints on
valid input. Thus the policy is competitive in its instrumented work unit,
up to the explicitly chosen input-dependent initial allowance. It is not
a bound of $12W$ when $W<L$.

The allowance formula is a scheduling choice, not a theorem that an AHT
run should finish in $O(kb)$ checkpoints. Its purpose is to let short runs
finish without paying for a second direction while retaining a geometric
restart bound when that initial allowance is exceeded. The proof does not
rely on the width selector being correct: it only changes which attempt
is tried first at each scale.

This does not prove constant-factor competitiveness in wall-clock time.
Checkpoint spacing is not uniform: input validation, sorting, heap
construction, identity filtering and some list operations can perform
multiple elementary operations between checkpoints. Big-integer division
and gcd also have operand-dependent cost. Each such block still has a
polynomial bound in the original pairing count and input bit length, and
the ordinary AHT run has a polynomial checkpoint count. Combining those
bounds with the displayed inequality preserves polynomial bit complexity.
The restart policy runs one search state at a time and retains only the
winning trace. None of these statements establishes a sharper general
bit-complexity bound for AHT or a bound for discovering a useful normal surface.

\subsection{Validation and performance consequences}

The frozen baseline is revision \code{cb26c81af}. Its interval producer,
interval checker, normal wrapper and outer checker are loaded separately,
with the old dependencies explicitly reconnected as in the preceding
section. The default compatibility audit covers all 1,275 maintained
surfaces, both coordinate-multiplicity settings and both boundary-classification
settings: 5,100 entire default records agree. All 5,100 race results have
the same topology, agree with the frozen component-level oracle and pass
independent replay. The coorientation shortcut and the geometric source
bindings remain unchanged.

The interval audit includes 500 random interval systems, 100 random
point-pair graphs, four star graphs and a fixed switching example. Both
merger thresholds give 1,210 literal component comparisons and checkpoint
bound checks. The three existing fixed-direction policies reproduce all
3,630 old complete records exactly. For every race, the audit checks the
actual caller checkpoint total, the sharper displayed bound and equality
of the winning certificate with its corresponding fixed-direction proof.

The switching example is a 72-point graph with 195 singleton pairings.
It deliberately exercises scheduling; ordinary graph union--find is a
more appropriate solver for such an expanded small graph. The race first
tries forward, exhausts five attempts and completes in reverse on attempt
six. The final allowance is 50,176 checkpoints and the successful run uses
50,174. It returns the same reverse certificate, after 175,619 total
attempt checkpoints and 196 startup checkpoints. Its 211 begun cycles
include 68 in the successful attempt. This supplies a concrete check of
switching and accounting, without presenting the synthetic graph as a
knot-recognition gain.

Focused tests check exact cycle caps including interrupted attempts,
caller cancellation, all three merger schedulers, integers exceeding
20,000 bits, empty universes, normal-surface composition and producer-free
replay. They also isolate the scheduler with deterministic checkpoint
streams so that second-direction completion and interrupted prefixes do
not depend solely on a particular geometric fixture.
All 16 focused tests and all 1,261 maintained tests pass.
The normal-surface audit exercises 2,972 restarts across its 5,100
configurations. In the default-multiplicity, unclassified corpus batch,
3,586 actual orbit queries use 4,264 attempts, including 678 restarts.
The race performs 4,724,655 attempt checkpoints and 65,555 startup checkpoints.
Its winning traces use 31,729 cycles, but its total begun-cycle count is
41,464 after including abandoned attempts. Forward uses 31,554 cycles.
This is why the public allowance must count the attempted work rather than
only the eventual certificate's search.

The race's returned corpus proofs contain 189,333 events and 13,289,274
serialized bytes. That is slightly smaller than the width-only result in
the previous section, but remains larger than forward's 182,925 events and
12,861,213 bytes. Competitive checkpoint scheduling is neither a proof-size
optimizer nor a guarantee that the first completed result is faster in
wall-clock time. The retained star and switching fixtures show the scale
of restart overhead directly:
\begin{center}
\small
\begin{tabular}{lrrrrr}
\toprule
Interval fixture & $W$ & attempt work $A$ & bound & attempts & total cycles\\
\midrule
Star 16 & 632 & 632 & 1537 & 1 & 16\\
Star 32 & 2544 & 6130 & 10755 & 3 & 62\\
Star 64 & 10208 & 34788 & 57349 & 5 & 166\\
Star 128 & 40896 & 169926 & 276487 & 7 & 403\\
Switching graph & 50174 & 175619 & 175621 & 6 & 211\\
\bottomrule
\end{tabular}
\end{center}


The isolated benchmark uses five randomized rounds, with a separate control
for each implementation. The first comparison uses frozen forward search
as its reference. Each timed surface call includes topology production,
independent replay, proof serialization and hashing; input construction
is outside the timer. Both arms retain coorientation and default multiplicity
reduction, and omit optional boundary classification.
\begin{center}
\small
\begin{tabular}{lrrrrr}
\toprule
Input & reference ms & race ms & ref./race & ref. A/A & race A/A\\
\midrule
Layered 16 & 10.833 & 9.297 & 1.167 & 0.981 & 0.987\\
Layered 64 & 110.830 & 58.769 & 1.884 & 0.978 & 0.973\\
Layered 256 & 1742.291 & 660.816 & 2.630 & 1.007 & 0.998\\
Empty & 0.363 & 0.436 & 0.833 & 1.000 & 1.036\\
M\"obius band & 0.697 & 0.829 & 0.846 & 0.983 & 1.012\\
Vertex-link disc & 1.460 & 1.668 & 0.885 & 0.983 & 1.017\\
Figure-eight surface & 9.407 & 21.093 & 0.447 & 1.003 & 0.984\\
Interior trefoil surface & 6.679 & 7.805 & 0.855 & 1.017 & 1.000\\
Entire 1,275-vector batch & 4885.100 & 6928.233 & 0.703 & 0.994 & 0.998\\
\bottomrule
\end{tabular}
\end{center}


The second comparison uses frozen width-selected search as its reference.
It isolates the extra cost of checkpoint accounting and restarts beyond
the same initial ordering preference. The final two rows time interval
production, independent interval replay and proof serialization directly;
they are scheduling stress cases, not normal-surface or knot pipelines.
\begin{center}
\small
\begin{tabular}{lrrrrr}
\toprule
Input & reference ms & race ms & ref./race & ref. A/A & race A/A\\
\midrule
Layered 256 & 500.632 & 650.542 & 0.762 & 1.005 & 1.001\\
Figure-eight surface & 10.847 & 20.928 & 0.518 & 1.001 & 0.995\\
Entire 1,275-vector batch & 4957.851 & 6916.716 & 0.716 & 0.995 & 0.994\\
Star 64 (interval kernel) & 4.229 & 19.827 & 0.216 & 0.996 & 0.997\\
Switching graph (kernel) & 24.510 & 106.682 & 0.230 & 1.021 & 0.993\\
\bottomrule
\end{tabular}
\end{center}


All 280 measured calls and 56 warm-ups finish with matching topology or
orbit-count digests. Proof and count records are stable within each arm.
Across the normal-surface cases there are 51,200 measured and 10,240 warm-up
surface analyses, counting every surface within the batch calls;
the interval rows add 40 measured and eight warm-up kernel calls.
Factors greater than one favor the race, and are medians of paired ratios.
Against forward, layered-256 improves by $2.630\times$, but the corpus
factor is $0.703$, about 42 percent slower. Against the width-only policy,
the corresponding factors are $0.762$ and $0.716$: the race is about
31 and 40 percent slower. The figure-eight counterexample is over twice
as slow as forward, and the two interval stress cases are over four times
as slow as their width-selected references. Control medians range from
$0.973$ to $1.036$, much smaller variation than these major regressions.
The comparison includes the scheduler's callback overhead and all abandoned
searches, while serializing only the winning proof. The optional policy
does not replace forward as the default. Avoiding repeated work through
resumable searches is a separate implementation candidate; these measurements
do not establish its performance.


The driver is \path{../fast/normal_orbit_research/race.py}, with
\code{audit} and \code{benchmark} modes, a required \code{--output} path,
and optional \code{--rounds}. Source pins, full switching certificates,
validation logs and timing records are retained as \path{data/orbit-race-*}.
The result adds a rigorously accounted adaptive search option. Promotion
to a default requires a favorable practical tradeoff, and a stronger
running-time guarantee would require finer accounting or resumable
implementations with bounded work between suspension points. The general
quasipolynomial recognition goal remains a separate open obligation for
this implementation.
